\documentclass[12pt,a4paper]{article}

% --- 基础环境配置 ---
\usepackage[UTF8,fontset=mac]{ctex}
\setCJKmainfont{Songti SC}
\setCJKsansfont{Heiti SC}
\usepackage{geometry}
\geometry{left=2cm,right=2cm,top=2.5cm,bottom=2.5cm}
\usepackage{amsmath, amssymb, amsfonts, bm}
\usepackage{graphicx}
\usepackage{xcolor}
\usepackage{titlesec}
\usepackage{hyperref}
\usepackage{cite}
\usepackage{abstract}
\usepackage{setspace} 
\usepackage{fancyhdr} 

% --- 样式定制 ---
\definecolor{scienceblue}{RGB}{0, 80, 158}
\hypersetup{colorlinks=true, linkcolor=scienceblue, citecolor=red, urlcolor=teal}
\onehalfspacing 

% 页眉页脚设置
\pagestyle{fancy}
\fancyhf{}
\fancyhead[L]{\small \color{gray} 刘晓华：凝聚态多体理论对人工智能的范式重构}
\fancyhead[R]{\small \color{gray} \thepage}

% 章节样式设置
\titleformat{\section}{\large\bfseries\color{scienceblue}}{}{0em}{\thesection.\quad}
\titleformat{\subsection}{\bfseries}{}{0em}{\thesubsection\quad}

\usepackage{gbt7714}

% --- 标题与个人信息 ---
\title{\Huge \textbf{从演生论到物理智能：\\凝聚态多体理论对人工智能的范式重构}}
\author{\Large \textbf{刘晓华} \\[0.5em]
\small 精细爆破全国重点实验室 \\ 
\small 数智爆破所}
\date{\today}

\begin{document}

\maketitle

\begin{abstract}
当前人工智能正处于从“关联驱动”向“原理驱动”转型的关键阶段。传统深度学习架构在大规模数据处理方面取得了显著进展，但在科学发现、逻辑稳健性和能效等方面暴露出明显不足。凝聚态物理（Condensed Matter Physics, CMP）作为研究多体系统演生行为（Emergent Behavior）的核心学科，为理解大规模参数系统的非线性动力学提供了自然的理论框架。本文尝试系统梳理重正化群（Renormalization Group, RG）在特征抽取与多尺度建模中的作用、张量网络（Tensor Networks, TN）在长程关联与维数约简中的优势，以及拓扑序（Topological Order）在构建鲁棒表征与抑制“幻觉”中的潜在价值。同时，结合 2023–2025 年 AI4Science 领域的若干代表性进展，讨论状态空间模型（SSM）、生成式扩散模型与随机重正化流之间的联系。最后，本文提出“全息语义重正化”“拓扑逻辑引擎”等若干值得进一步探索的交叉方向，希望为面向物理规律的通用智能系统建设提供相对统一的理论视角与研究议程。
\end{abstract}

\newpage
\renewcommand{\contentsname}{目 录}
\tableofcontents
\newpage

% --- 第一章：物理智能的哲学基础 ---
\section{引言：More is Different 与 AGI 的演生}

过去十年间，人工智能（AI）经历了从“小模型+手工特征”向“大模型+大数据+大算力”的快速演化。随着参数规模从亿级迈向万亿级，研究者逐渐意识到，当前大模型表现出的若干“能力跃迁”现象，很难仅用传统的“统计拟合”范式加以解释。这一局面，与凝聚态物理在研究大规模电子、原子和自旋体系时面临的“多体问题”高度相似：单个粒子的微观定律是清楚的，但宏观涌现的结构和功能，往往并非简单叠加的结果。

凝聚态物理以多体相互作用、相变和演生为中心问题，提供了一套从微观到宏观、从局域到全局的“尺度组织”视角。将这一视角引入 AI，不仅有助于理解当前大模型的若干“怪现象”，也有可能为下一代智能系统的设计提供更加克制、可验证的理论依据。

\subsection{演生论（Emergence）：连接电子气与大模型的桥梁}

1972 年，P.W. Anderson 在经典论文《More is Different》中\cite{anderson1972more} 对还原论（Reductionism）提出了著名的批评：当系统复杂性达到一定程度后，宏观层面会出现新的规律，这些规律并不能从微观方程的形式直接推演出来。在超导、超流、自旋玻璃等体系中，单电子或单自旋的动力学方程并不足以解释零电阻、长程有序或复杂能量景观等性质；这些性质是在热力学极限 $N \to \infty$ 和强相互作用背景下“演生”出来的。

大语言模型（LLM）在缩放到一定规模之后表现出的少样本学习、思维链推理等能力，同样呈现出类似“相变式”的跃迁特征。与其简单地把这些现象理解为“拟合能力增强”，不如借用凝聚态物理的语言，将其视为参数空间中的一种\textbf{新相}：模型从“无序随机映射”演化为具有内在结构和长程关联的“有效低能相”。

从这一角度看，我们可以将智能视作参数空间中的一种\textbf{低能有效相（Effective Low-energy Phase）}。就像有效场论（Effective Field Theory, EFT）允许在不显式追踪每个粒子轨迹的前提下描述流体性质一样，对大模型而言，单个神经元的具体权重不再重要，更重要的是在大尺度上形成的对称性、守恒量与相结构。AI 研究由此可以从“调参经验学”转向“相结构与演生机制”的研究。

\subsection{从统计物理到神经计算的诺奖脉络}

统计物理与智能计算的交叉，并非近年的偶然现象，而是有着清晰的历史脉络。2021 年 Giorgio Parisi 因自旋玻璃和复杂系统的工作获诺贝尔物理学奖\cite{parisi1980order}，而 2024 年 John Hopfield 与 Geoffrey Hinton 的获奖，则进一步确认了“统计物理—神经网络”这一方向的重要性。

\subsubsection{Hopfield 模型与能量最小化}

Hopfield 在 1982 年的工作\cite{hopfield1982neural} 中提出，将联想记忆建模为一个伊辛型多体系统的能量最小化过程。其能量函数
\begin{equation}
    E = -\frac{1}{2}\sum_{i \neq j} w_{ij} s_i s_j
\end{equation}
刻画了“模式”作为能量景观的极小值。学习即塑造能量景观，记忆即在能量景观中的吸引子间演化。这个视角为“智能是一种统计力学现象”提供了最直接的物理表达。

\subsubsection{Parisi 理论与损失景观}

Parisi 在自旋玻璃中的复件对称性破缺（RSB）理论，为理解复杂损失景观提供了成熟工具。现代大模型的损失函数空间高维、非凸且充满局部低谷，Biroli 等人的工作\cite{biroli2024energy} 表明，可将其视作一种类似 $p$-自旋玻璃的能量景观，训练过程则类似低温退火。诸如梯度消失、训练停滞、泛化能力差等问题，可以在遍历性破缺、亚稳态和能垒结构的框架下重新审视。

\subsubsection{“物理智能”的工作假说}

从上述历史脉络出发，一个愈发清晰的判断是：走向更高层次智能（包括 AGI）的道路，很难仅靠算力堆叠。更有前景的路径，是在模型结构、训练目标和评估体系中，系统性地引入物理规律，特别是与演生、相变、拓扑和非平衡相关的结构性约束。本文后续各章，将围绕这一工作假说展开更具体的讨论。

% --- 第二章：重正化群——学习即多尺度信息截断 ---
\section{重正化群（RG）：深度学习的多尺度过滤本质}

在凝聚态物理中，重正化群（RG）既是一套计算方法，也是一种处理多尺度问题的基本思想：通过逐步“积掉”高能自由度，保留对低能行为起决定作用的有效变量。类比到深度学习，多层网络天然具有“逐层抽取特征”的结构，这与 RG 的多尺度过滤十分接近。

\subsection{变分重正化群与层级特征映射}

Mehta 与 Schwab 在 2014 年的工作\cite{mehta2014exact} 中，给出了受限玻尔兹曼机（RBM）与变分重正化群之间的严格对应关系。

\subsubsection{从微观状态到有效哈密顿量}

对于一个格点 Ising 模型，其配分函数 $Z = \sum_{\{s\}} e^{-H(s)}$ 可通过块自旋变换映射到粗粒化变量 $\{h\}$。在 RBM 中，可见层 $\{s\}$ 与隐藏层 $\{h\}$ 之间的条件分布
\begin{equation}
    P(h) = \sum_{\{s\}} P(h|s) P(s) \propto \sum_{\{s\}} e^{H_{RBM}(s, h)}
\end{equation}
可以理解为对高频自由度的积分，得到新的“有效哈密顿量”。隐藏层起到的作用，正是一次“紫外截断”：抛弃局部噪声，保留对后续表示仍有影响的关联结构\cite{lin2017why}。沿层数方向演化，可以视作在“参数—尺度空间”上的一条 RG 流。

\subsection{信息瓶颈：学习、压缩与泛化}

Tishby 的信息瓶颈（IB）理论\cite{tishby2015deep} 为 RG 与深度学习的联系提供了信息论视角。网络训练往往经历一个快速拟合阶段和一个缓慢压缩阶段：前者主要捕捉与标签相关的模式，后者则在保持预测性能的前提下压缩无关信息。IB 模型指出，好的表征应当在“保留任务相关信息”和“压缩输入冗余”之间取得平衡，这一平衡可视作 RG 流找到的某个“不动点”（Fixed Point）\cite{koch2018mutual}。

对于 AI4Science，特别是在物理实验数据往往噪声较大、样本稀缺的背景下，这一思想具有实际意义：与其盲目扩展模型容量，不如通过引入物理先验，引导网络在 RG 意义上收敛到“合物理”的表示上。

\subsection{扩散模型与随机重正化流}

扩散模型近年的成功，为 RG 思想在生成建模中的应用提供了新舞台。近期工作\cite{hu2024learning, chen2025diffusion} 指出，扩散过程的“加噪—去噪”可以理解为一条“从无序到有序”的随机 RG 轨迹。

\subsubsection{逆向过程与 RG 流}

在正向扩散中，图像逐渐被噪声淹没，可视为向高温无序定点流动；在逆向生成中，模型通过得分函数引导系统从高温无序回到“结构良好”的低温相。借助类似 Callan–Symanzik 方程的工具，
\begin{equation}
    \left[ \frac{\partial}{\partial \ln \Lambda} + \beta(g) \frac{\partial}{\partial g} + \gamma \right] \Gamma^{(n)} = 0
\end{equation}
可以尝试将网络参数 $g$ 的演化视作 RG 流，从而讨论不同尺度上纹理与形状、局部与整体之间的权衡。

\subsection{对工程科学建模的启示}

在精细爆破等多尺度工程问题中，微观裂纹萌生、介质损伤和宏观波场演化跨越多个数量级。若能在网络结构中显式引入 RG 逻辑，则有望构造一种“尺度自适应”的模型：在微观层面使用物理约束强的 PINNs\cite{raissi2019physics} 捕捉局部行为，在宏观层面通过重正化式压缩将微观细节折算为有效本构参数，从而提高预测的稳定性和可解释性。

% --- 第三章：张量网络——信息几何与纠缠熵约束 ---
\section{张量网络：纠缠熵约束下的信息几何与架构革新}

量子多体问题的核心难点在于希尔伯特空间维度随粒子数指数增长。张量网络通过在结构上引入局域关联和有限纠缠假设，将高维问题嵌入到一个低维流形上，从而巧妙地绕开了“维数灾难”。这一思路对于理解和改造深度学习架构具有直接启发。

\subsection{纠缠熵面积定则与神经网络稀疏结构}

面积定则（Area Law）\cite{eisert2010colloquium} 告诉我们，自然界中许多基态的纠缠熵与区域边界面积而非体积成正比。这意味着“物理上真实的量子态”只占据希尔伯特空间的极小一部分。

类比于此，自然图像、语音和文本的有效自由度远小于像素或符号的组合空间。Stoudenmire 和 Schwab\cite{stoudenmire2016supervised} 利用矩阵乘积态（MPS）构建监督学习模型，说明在适当的键维（Bond Dimension $D$）下，可以用远少于全连接网络的参数实现可比性能。这从物理角度解释了为什么大多数任务并不需要“全连接注意力”，也为极限压缩、边缘部署提供了基础。

\subsection{SSM 与张量链：线性动力学的物理根基}

2023–2024 年出现的 Mamba 等状态空间模型\cite{gu2023mamba}，本质上采用了更接近连续时间动力学的建模方式。最近一些工作\cite{ali2024tensor} 将其与非均匀 MPS 链联系起来：隐藏状态的递推可以看作沿张量链方向的传递矩阵演化，
\begin{equation}
    W_{s_1 s_2 \dots s_n} = \text{Tr}(A_1^{s_1} A_2^{s_2} \dots A_n^{s_n})
\end{equation}
这种结构天然具有 $O(N)$ 的时空复杂度，有望在长序列场景中取代 Transformer 的二次复杂度，并与物理中的格林函数、传播子形式建立更紧密的联系。

\subsection{MERA 与全息语义表征}

多尺度纠缠重正化（MERA）\cite{vidal2008class} 在处理临界系统时表现突出，其树状结构与 AdS/CFT 全息对偶之间的对应早已广为人知。在语言建模中，可以考虑用 MERA 结构构造一种“自带尺度”的语义表示：底层节点对应局部搭配与句法，中层节点对应句子级结构，顶层节点则对应篇章级意图。相比扁平的自注意机制，全息式结构更方便引入跨尺度一致性约束。

\subsection{高维场问题中的张量分解}

在爆破、地震和复杂材料模拟中，高维场（应力、速度、温度等）网格数据往往达到 GB 级别。Tensor Train 等张量分解方法\cite{novikov2015tensor} 可将这些高维数组表示为一系列低阶张量的乘积，在保证精度的前提下实现大幅压缩。一方面可减轻存储压力，另一方面也便于在低秩流形上进行物理约束下的优化，为实时数值模拟和“数智孪生”构建提供基础。

% --- 第四章：统计力学视角下的涌现与相变 ---
\section{统计力学视角下的涌现与相变：从自旋玻璃到能力门槛}

如果说 RG 和张量网络解决的是“多尺度”和“高维”的问题，那么自旋玻璃与统计力学则为理解“能力突变”和“训练动力学”提供了语言。

\subsection{Hopfield 模型：联想记忆与能量景观}

前文已提及 Hopfield 模型\cite{hopfield1982neural} 将记忆模式视作能量景观中的极小值。现代大模型虽然结构更复杂，但“能量景观—吸引子—动力学”的主线仍然适用：训练过程中，参数沿着高维损失景观的斜坡缓慢下滑；推理过程中，激活态在隐空间的不同吸引盆之间演化。如何塑造有利于泛化的能量景观，如何避免“坏盆地”，本质上是统计力学问题。

\subsection{自旋玻璃、RSB 与损失景观}

Parisi 的 RSB 理论\cite{parisi1980order} 揭示了多谷景观、亚稳态和复杂能量层次结构的普遍性。Biroli 等人的工作\cite{biroli2024energy} 指出，可用类似 $p$-自旋模型刻画深度网络的训练。在高“温度”（如大学习率、强噪声）下，系统处于无序相；随着“降温”（降低学习率、增加训练轮数），系统可能经历一次或多次“相变”，进入具有明显结构的有序相。

\subsection{大模型涌现能力的“相变”解释}

近年来关于“涌现能力是否真实存在”的争论\cite{schaeffer2024are}，可以在相变框架下重新审视。如果将模型能力的某一指标视作“序参量”，将参数规模、数据量或训练步数视作“控制参数”，则在某些区域确实可以观察到类似二级相变的非解析性：能力曲线出现拐点，响应函数（例如能力对参数的导数）出现峰值甚至接近发散。无论这种现象是否构成严格数学意义上的相变，其物理类比都对理解和控制模型缩放具有启发意义。

\subsection{自组织临界性与“混沌边缘”}

Bak 等人提出的自组织临界性（SOC）\cite{bak1987self} 认为，许多复杂系统天然演化到临界点附近，其特征是幂律分布和长程关联。神经元放电、脑电活动等已有大量研究表明，生物大脑常常在“临界区”附近运行\cite{chialvo2010emergent}。近期关于大模型内部激活统计的工作\cite{zimmern2024large} 也呈现出类似特征。这一现象提示我们，在网络结构与训练策略设计中，有意识地引导系统靠近“混沌边缘”，可能有助于提升表达能力和适应性。

% --- 第五章：拓扑序与逻辑鲁棒性 ---
\section{拓扑序与逻辑鲁棒性：抑制幻觉的物理思路}

当前大模型在推理和事实性方面的缺陷，很大程度上源于缺乏“全局保护机制”：局部错误容易在层间演化中被放大。拓扑序为构建具有全局稳定性的表示提供了一个值得参考的方向。

\subsection{贝里相位与持续学习中的表征几何}

在持续学习或在线微调中，参数和表征在训练流形上不断移动。若忽略这一几何结构，容易出现灾难性遗忘和逻辑漂移。借鉴贝里相位（Berry Phase）\cite{berry1984quantal} 概念，可将表征演化视为参数空间上的平行移动：
\begin{equation}
    \boldsymbol{\mathcal{A}}(\mathbf{w}) = i \langle \Psi(\mathbf{w}) | \nabla_{\mathbf{w}} | \Psi(\mathbf{w}) \rangle
\end{equation}
在此基础上引入几何正则项，有望在训练过程中约束表征的“路径”，从而减小逻辑偏移。

\subsection{拓扑不变量与鲁棒推理}

拓扑绝缘体的一个重要特征，是体态存在能隙而边缘态具有鲁棒性，这由陈数等拓扑不变量保证。类似地，若能在神经表示或逻辑结构中构造某种“拓扑不变量”，就有可能实现对局部噪声的全局抑制。

近期综述\cite{battiloro2024topological} 汇总了“拓扑神经网络”的若干尝试，包括在特征空间中引入具有非平凡拓扑的流形结构等。Fujimoto 等工作\cite{fujimoto2024topological} 进一步提出利用拓扑能隙概念刻画表示的稳定性。这一方向尚处早期，但在安全可信 AI、对抗鲁棒等场景中值得重点关注。

\subsection{准粒子编织与离散逻辑}

拓扑量子计算以非阿贝尔任意子（Anyons）的编织实现逻辑门，其优势在于运算结果仅依赖于轨迹的拓扑类型而非具体细节。若能在离散逻辑推理中引入类似“拓扑编织”的抽象结构，使推理序列对局部扰动不敏感，则有望构造一种“拓扑保护”的逻辑引擎，从根本上提高推理的可复现性和稳定性。

\subsection{AI4Science 中的拓扑奇异点}

在爆破与复杂材料行为中，位错、裂纹尖端、涡旋等拓扑缺陷往往决定宏观响应。传统数值方法在处理这些奇异点时代价高昂且稳定性较差。借助拓扑数据分析（TDA）\cite{carlsson2009topology} 和拓扑场论引导的网络结构，有望更准确地捕捉这些关键结构，进而提升模型对爆破过程分叉和失效模式的预测能力。

% --- 第六章：未来研究方向 ---
\section{未来研究方向：迈向“物理智能”的若干可能路径}

综合前述讨论，从凝聚态物理出发理解和重构智能系统，并非简单的比喻，而是可以形成一套相对系统的研究框架。以下列出若干我认为值得中长期投入的方向。

\subsection{全息语义学：基于 MERA/AdS-CFT 的语言架构}

现有 Transformer 架构在本质上仍是欧氏空间上的扁平注意力。若引入 MERA 结构\cite{vidal2008class}，并将其与 AdS/CFT 全息对偶类比，有机会构造一种自带尺度的语义表示体系：高层节点对应全局意图，低层节点对应局部实现。通过在训练中显式约束“体—边界”一致性，有可能缓解长文本生成中前后矛盾的问题。

\subsection{拓扑逻辑引擎：受拓扑保护的推理}

受拓扑量子计算启发，可以探索将推理过程离散化为“拓扑操作序列”，并为之定义某种“逻辑能隙”。只要干扰不超过该能隙，推理结果即保持不变。这类工作一旦取得实质性进展，将有望在形式化证明、复杂工程安全评估等对可靠性要求极高的场景中发挥作用。

\subsection{算力与能效的重正化：自适应激活与剪枝}

现实应用中，大模型在绝大多数请求下并不需要全容量参与。若能构造一套类似 RG 的算子流方程，
\begin{equation}
    \frac{d g(\ell)}{d \ell} = \beta(g)
\end{equation}
根据任务难度和输入结构，自适应激活子网络并进行“动态剪枝”，则有望实现“按需用脑”，大幅提升能效。这与生物大脑在不同任务负载下激活不同网络的经验观察相一致。

\subsection{非平衡态演生智能：从静态映射到耗散结构}

现有大模型更像准静态映射，而真实智能体通常处于非平衡态。将网络视作一种活性物质系统\cite{vicsek2012collective}，研究信息流在能量输入和耗散背景下如何形成稳定结构，可能有助于理解“具身智能”和“自主性”。在爆破等极端条件下材料的快速相变问题上，此类框架也有望提供新的数值方法。

% --- 第七章：工程实现挑战与路线图 ---
\section{工程实现挑战与 AI4Science 实践路径}

从理论走向工程落地，还需要正视一系列现实约束。

\subsection{大规模张量收缩与硬件瓶颈}

张量网络在理论上优雅，但在工程上面临收缩顺序优化和中间张量爆炸的问题。现有 GPU/TPU 更偏向矩阵乘法，对复杂多指标收缩支持有限。需要结合 cuTensor 等库，发展面向张量网络的编译与调度策略，以便在大规模科学模拟中真正发挥 TN 的优势。

\subsection{小样本与物理一致性}

在精细爆破、材料失效等领域，高质量数据难以大量获取。单纯依赖数据驱动容易过拟合。更现实的路径，是在损失函数中引入物理守恒、对称性和尺度律等约束，将“物理先验”与“数据驱动”结合，使模型在小样本条件下仍保持物理可接受性。

\subsection{硬件架构的适配与演进}

冯·诺依曼架构天生存在“存算分离”的瓶颈。面向下一代 AI4Science 任务，值得关注忆阻器、神经形态芯片、自旋电子学等新型硬件，它们在物理层面更接近连续场演化，有潜力与张量网络、能量模型等结构自然耦合，降低复杂科学计算的能耗门槛。

% --- 第八章：结语 ---
\section{结论：从“炼丹”走向“第一性原理”的演生智能}

本文尝试从凝聚态物理的视角，对当前 AI 与 AI4Science 的若干关键问题作一初步梳理。 Anderson 的演生论、重正化群、张量网络、拓扑序、自旋玻璃与自组织临界性等概念，为理解大模型的能力涌现、有效表示和鲁棒性问题提供了一套相对统一的语言。

对于精细爆破全国重点实验室和数智爆破所而言，真正有价值的不是“再造一个 ChatGPT”，而是建设一个能够理解和尊重物理规律的“数字伙伴”：在多尺度、多物理场、多不确定性的复杂环境中，帮助工程师和科学家更安全、更高效地做出决策。我个人相信，将凝聚态物理与智能计算有机结合，不仅可以推动 AI4Science 的具体应用，也有机会反哺基础物理本身，对“复杂系统如何在多体相互作用中演生出结构与功能”这一根本问题，提供新的视角。

未来的智能，或许不应被简单地理解为更大的模型和更快的芯片，而应被理解为物理定律、信息论与工程实践在高维空间中的一次持续协同演化。

% --- 参考文献 ---
\newpage
\bibliographystyle{gbt7714-numerical}
\bibliography{refs}
 


\end{document}
